\chapter{From Angular Momentum to Spin}

What is formally announced as a lecture on angular momentum quickly narrows to the real target: spin. We begin from ordinary rotation because that is the natural classical entry point, but the point of the chapter is to see how a very spare piece of quantum algebra produces the characteristic facts about spin. The order matters. First we separate orbital angular momentum from spin, then we ask when a rotating state should count as the same object and when it should count as a different one, and only after that do we turn to the mathematics that makes the subject so striking. These notes follow Leonard Susskind's lecture, with curation by LazyingArt LLC.

\section{From Ordinary Rotation to Spin}

Angular momentum is first introduced in the most elementary way: it is a vector quantity, and its direction is fixed by the right-hand rule. If a body rotates about an axis, the angular momentum points along that axis, but the sign convention is not proved; it is defined. One wraps the fingers in the sense of rotation and the thumb gives the direction of the angular-momentum vector.

At this stage the lecture insists on a distinction that will govern everything that follows. A composite object may have angular momentum simply because its center of mass moves around some axis. That is orbital angular momentum. But the angular momentum we are after is the angular momentum that remains when we go to the rest frame of the center of mass. In that frame the total ordinary momentum vanishes. If there is still angular momentum left over, that residual angular momentum is what we call spin.

This distinction is already visible in a very simple classical picture. A single particle moving on a circle has orbital angular momentum because it has momentum and a lever arm relative to the origin. But if two particles sit at opposite points of a diameter and rotate with the same sense, then the center of mass can be at rest while the angular momentum is not zero. That is the lecture's first concrete model for spin.

The next question is not yet algebraic. Suppose an object begins with no angular momentum and we excite it into a state with nonzero angular momentum. Is it still the same object, now rotating, or should we regard it as a different object? Susskind treats that as partly a question of language. The more physical question is the energy gap to the first rotationally excited state. In classical mechanics we can spin an object up by an arbitrarily small amount, so the cost can be made arbitrarily small. In quantum mechanics the rotational states are discrete, and that makes the first step meaningful. If the first step costs very little energy, we tend to recognize the rotating system as the same object. If it costs an enormous amount, the excitation may be better thought of as a different object altogether.

This becomes especially sharp for an electron. If we forget, for a moment, that the electron already has intrinsic spin, we may still ask whether a tiny object can be ``set rotating'' in the ordinary sense. The lecture's answer is cautious but clear: for very small objects the gap to the next rotational state may be so large that we would not describe the result as the same particle with more spin. We would describe it as a different object. In this sense the spin of an elementary particle behaves as a fixed label.

\subsection{Question \& Answer}

Why does a smaller object cost more energy to spin?

The answer begins classically, with the moment of inertia. For a body of mass \(M\) and size \(R\), the moment of inertia is of order
\begin{equation}
I \sim MR^2,
\end{equation}
up to shape-dependent numerical factors. The rotational kinetic energy is then of order
\begin{equation}
E_{\mathrm{rot}} \sim \frac{L^2}{2I}.
\end{equation}
So if the mass is held fixed and the size \(R\) is made smaller, then \(I\) becomes smaller, and because \(I\) sits in the denominator the same angular momentum costs more energy. Quantum mechanics adds only one crucial ingredient: the allowed angular momenta come in discrete steps. The first nonzero rotational excitation of a very small object can therefore be extremely expensive.

This is the physical reason the lecture treats the intrinsic spin of an elementary particle as a fixed property. One does not ordinarily speak of electrons with many different total spin magnitudes. The gap to the next such state is so large that one would no longer be talking about the same familiar particle.

\section{Classical Angular Momentum As a Cross Product}

Once the physical motivation is in place, the lecture turns to the classical formula. For a single particle orbiting an origin, the angular momentum is momentum times distance from the origin. In vector language, the position is described by the radius vector \(\vec r\), and the momentum by \(\vec p\). The angular momentum must itself be a vector built from these two vectors. The appropriate rule is the cross product:
\begin{equation}
\vec L=\vec r\times \vec p.
\end{equation}

The board at this point is visually faint, but the geometric setup and the transcript both support this standard formula.

\begin{figure}[h]
\centering
\includegraphics[width=0.72\textwidth]{lecture_07_figure_02.png}
\par
\begin{tikzpicture}[scale=1.0]
\draw[->] (-0.6,0) -- (4.8,0) node[right] {$x$};
\draw[->] (0,-0.6) -- (0,3.8) node[above] {$y$};
\fill (0,0) circle (1.3pt);
\coordinate (P) at (3.2,2.3);
\draw[->,thick] (0,0) -- (P) node[midway, above left] {$\vec r$};
\draw[dashed] (P) -- (3.2,0);
\draw[dashed] (P) -- (0,2.3);
\draw[->,thick] (P) -- ++(-1.0,1.25) node[above] {$\vec p$};
\end{tikzpicture}
\caption{Angular momentum from position and momentum vectors. The board image records the geometric setup; the sketch below gives a clean reconstruction of the move to \(\vec L=\vec r\times\vec p\).}
\end{figure}

If we now resolve the cross product into components, we obtain
\begin{equation}
L_x=yp_z-zp_y,\qquad
L_y=zp_x-xp_z,\qquad
L_z=xp_y-yp_x.
\end{equation}
The lecture emphasizes a practical mnemonic rather than a formal determinant. One remembers one component, say
\[
L_x=yp_z-zp_y,
\]
and then cycles \(x\to y\to z\to x\) to generate the others. The point of the cycling rule is not merely memory; it keeps the signs straight.

These formulas describe orbital angular momentum for a point particle near an origin. Spin for an ordinary composite system is obtained by adding angular momenta vectorially over constituents. If the total ordinary momentum vanishes in the center-of-mass frame while the total angular momentum does not, then the system has spin.

\section{Operators And The Angular-Momentum Algebra}

The lecture now makes the quantum-mechanical turn. Observable quantities are represented by operators, and for present purposes we do not need the full machinery of wavefunctions or representations. We only need the commutation relations of position and momentum. The stable canonical relations are
\begin{equation}
[x_i,x_j]=0,\qquad [p_i,p_j]=0,\qquad [x_i,p_j]=i\hbar\,\delta_{ij}.
\end{equation}
The transcript briefly garbles the discussion of coordinate commutators, but the derivation immediately proceeds with these standard relations, and these are the ones the notes should preserve.

The physical meaning is familiar. If two operators commute, they can be simultaneously measured. If they do not commute, they cannot be simultaneously specified with arbitrary precision. In particular, different position components commute with one another, different momentum components commute with one another, but a position component and the corresponding momentum component do not.

\paragraph{A worked commutator.}

The lecture explicitly computes \([L_x,L_y]\). Using
\[
L_x=yp_z-zp_y,\qquad L_y=zp_x-xp_z,
\]
we write
\begin{align}
[L_x,L_y]
&=[yp_z-zp_y,\;zp_x-xp_z] \\
&=[yp_z,\;zp_x]+[-zp_y,\;-xp_z].
\end{align}
The cross terms vanish because their factors commute completely. In the first surviving term the only nontrivial commutator is \( [p_z,z]=-i\hbar \), so
\begin{equation}
[yp_z,\;zp_x]=y[p_z,z]p_x=-i\hbar\,yp_x.
\end{equation}
In the second surviving term the only nontrivial commutator is \( [z,p_z]=i\hbar \), so
\begin{equation}
[-zp_y,\;-xp_z]=[zp_y,\;xp_z]=x[z,p_z]p_y=i\hbar\,xp_y.
\end{equation}
Adding the two pieces,
\begin{equation}
[L_x,L_y]=i\hbar(xp_y-yp_x)=i\hbar L_z.
\end{equation}
By cyclic permutation we obtain the full algebra,
\begin{equation}
[L_x,L_y]=i\hbar L_z,\qquad
[L_y,L_z]=i\hbar L_x,\qquad
[L_z,L_x]=i\hbar L_y,
\end{equation}
or, more compactly,
\begin{equation}
[L_i,L_j]=i\hbar\,\epsilon_{ijk}L_k.
\end{equation}

This is one of the lecture's central recap points. Once these commutators have been found, the algebra closes on the \(L_i\) themselves. We no longer need to carry the original \(x_i\) and \(p_i\) around. The subject has become abstract in the precise sense Susskind emphasizes: the algebra no longer remembers where it came from. It is also rotationally invariant. If we choose any other orthogonal axes, the commutators take the same form.

\section{Ladder Operators And The Quantization Axis}

At this point the lecture chooses a working axis. There is nothing physically special about the \(z\)-axis, but there is also nothing wrong with focusing on one axis at a time. In fact we must do so, because the components of angular momentum do not commute. At most one component can be sharply specified at once. This chosen axis is the quantization axis.

The new operators are
\begin{equation}
L_\pm=L_x\pm iL_y.
\end{equation}
On the blackboard the lecture temporarily sets \(\hbar=1\), but we restore \(\hbar\) in the final formulas. If
\begin{equation}
L_z\lvert m\rangle=m\hbar\,\lvert m\rangle,
\end{equation}
then the commutators may be written in the standard form
\begin{equation}
[L_z,L_\pm]=\pm \hbar L_\pm.
\end{equation}
The lecture often writes the equivalent reversed-order relations
\[
[L_\pm,L_z]=\mp \hbar L_\pm,
\]
and the two conventions are of course the same statement.

Now take a state \(\lvert m\rangle\) of definite \(L_z\). Then
\begin{align}
L_z\bigl(L_+\lvert m\rangle\bigr)
&=\bigl(L_+L_z+[L_z,L_+]\bigr)\lvert m\rangle \\
&=\bigl(L_+L_z+\hbar L_+\bigr)\lvert m\rangle \\
&=(m+1)\hbar\,\bigl(L_+\lvert m\rangle\bigr).
\end{align}
So \(L_+\lvert m\rangle\) is again an eigenstate of \(L_z\), now with eigenvalue \((m+1)\hbar\). Similarly,
\begin{equation}
L_-\lvert m\rangle \propto \lvert m-1\rangle.
\end{equation}
This is the first genuinely magical step in the lecture. The algebra alone tells us that the possible \(m\)-values must be spaced in integer steps.

\subsection{Question \& Answer}

If the ladder always raises or lowers by one step, how can the spectrum ever stop?

The raising argument proves only that if \(L_+\lvert m\rangle\) is nonzero, then it has eigenvalue \(m+1\). It does not prove that the vector must be nonzero. For a fixed particle the total spin magnitude is not supposed to grow without bound. So the only way the ladder can terminate is the same way termination occurs for the annihilation operator of the harmonic oscillator: at some point the result must vanish. Thus there must be a highest state satisfying
\begin{equation}
L_+\lvert m_{\max}\rangle=0.
\end{equation}
That condition is not an extra decorative fact. It is exactly what allows a finite multiplet to exist.

\section{Finite Multiplets, Highest Weight, And \texorpdfstring{$L^2$}{L2}}

Once the possibility of a highest state is admitted, the spectrum becomes sharply constrained. By repeated action of \(L_-\) we descend in unit steps. By rotational invariance the spectrum must be symmetric: if one can orient the angular momentum maximally upward, one must also be able to orient it maximally downward with the same magnitude. Therefore
\begin{equation}
m_{\min}=-m_{\max}.
\end{equation}
So the allowed values are
\begin{equation}
m=-m_{\max},-m_{\max}+1,\dots,m_{\max}.
\end{equation}
The total number of states is therefore
\begin{equation}
2m_{\max}+1.
\end{equation}

\begin{figure}[h]
\centering
\includegraphics[width=0.70\textwidth]{lecture_07_figure_03.png}
\par
\begin{tikzpicture}[scale=1.0]
\draw[thick] (1.6,-2.2) -- (1.6,2.2);
\draw (1.35,1.1) -- (1.85,1.1);
\draw (1.35,0) -- (1.85,0);
\draw (1.35,-1.1) -- (1.85,-1.1);
\draw (1.42,1.8) -- (1.78,1.8);
\draw (1.42,-1.8) -- (1.78,-1.8);
\node[right] at (1.95,1.1) {$m=1$};
\node[right] at (1.95,0) {$m=0$};
\node[right] at (1.95,-1.1) {$m=-1$};
\node[left] at (0.9,1.8) {$L_+\lvert m_{\max}\rangle=0$};
\end{tikzpicture}
\caption{Highest-weight condition and the \(m\)-level ladder. The board image records the cutoff condition together with the vertical spectrum; the diagram below gives a clean reconstruction.}
\end{figure}

At this stage the lecture separates the possibilities into two families. In one family the spectrum contains \(m=0\), so all \(m\)-values are integers. In the other family \(0\) is absent, and the spectrum consists entirely of half-integers:
\[
m=\pm \frac12,\ \pm \frac32,\ \dots,\ \pm m_{\max}.
\]
Only after this structure is established do we make the standard notational identification
\begin{equation}
s:=m_{\max}.
\end{equation}
The particle is then said to have spin \(s\). Thus \(s=0,1,2,\dots\) gives the integer family, while \(s=\tfrac12,\tfrac32,\dots\) gives the half-integer family.

The next invariant is \(L^2\). Naively, one might expect the top state to have \(L^2=s^2\hbar^2\), because the spin is pointing entirely along the \(z\)-axis. The lecture pauses to explain why this is too naive. If \(L_z\) is sharp, then \(L_x\) and \(L_y\) cannot both vanish sharply, because the components do not commute. There is an irreducible quantum correction.

The board algebra at this point is hesitant, so it is better to preserve the lecture's intent by writing the standard identity in its canonical form:
\begin{equation}
L^2=L_-L_+ + L_z^2 + \hbar L_z.
\end{equation}
Equivalently, on the board with \(\hbar=1\),
\[
L^2=L_-L_+ + L_z^2 + L_z.
\]

\begin{figure}[h]
\centering
\includegraphics[width=0.76\textwidth]{lecture_07_figure_04.png}
\caption{\(L^2\) acting on the highest-weight state. The board records the structure of the highest-weight calculation, but the final equations are best written in the standard restored-\(\hbar\) form in the text.}
\end{figure}

Now let this operator act on the highest-weight state:
\begin{align}
L^2\lvert s,s\rangle
&=\bigl(L_-L_+ + L_z^2 + \hbar L_z\bigr)\lvert s,s\rangle \\
&=L_-L_+\lvert s,s\rangle + s^2\hbar^2\lvert s,s\rangle + s\hbar^2\lvert s,s\rangle.
\end{align}
But \(L_+\lvert s,s\rangle=0\), so the first term vanishes. Therefore
\begin{equation}
L^2\lvert s,s\rangle=s(s+1)\hbar^2\,\lvert s,s\rangle.
\end{equation}
This is the characteristic quantum correction. Classically one would expect \(s^2\), but quantum mechanically one gets \(s(s+1)\). The extra \(+\,s\) comes from noncommutativity.

A further important fact is
\begin{equation}
[L^2,L_i]=0,\qquad i=x,y,z.
\end{equation}
So although one cannot simultaneously specify two components of \(\vec L\), one can specify \(L^2\) together with one chosen component, for instance \(L_z\). For a fixed spin \(s\), the natural basis of spin space is therefore
\[
\lvert s,m\rangle,\qquad m=-s,-s+1,\dots,s.
\]
Every state in the multiplet has the same \(L^2\) eigenvalue \(s(s+1)\hbar^2\), and the ladder operators merely move us from one \(m\)-value to another within the same irreducible set of states.

\section{Spin-\texorpdfstring{$\tfrac12$}{1/2} States And The Meaning Of Orientation}

The lecture next asks a very natural question. If a spin-\(\tfrac12\) particle has only the two \(z\)-axis values \(+\hbar/2\) and \(-\hbar/2\), what becomes of directions in between? This is where the language of states becomes unavoidable.

For spin \(\tfrac12\), let
\[
\lvert+\rangle=\left\lvert \frac12,\frac12\right\rangle,\qquad
\lvert-\rangle=\left\lvert \frac12,-\frac12\right\rangle.
\]
These are the basis states of definite \(L_z\). A general spin state is
\begin{equation}
\lvert\psi\rangle=\alpha\lvert+\rangle+\beta\lvert-\rangle,
\qquad
|\alpha|^2+|\beta|^2=1.
\end{equation}
Here \(|\alpha|^2\) and \(|\beta|^2\) are the probabilities for measuring \(+\hbar/2\) and \(-\hbar/2\) along the \(z\)-axis. If we multiply both coefficients by the same overall phase \(e^{i\theta}\), no physical prediction changes. So the state is really determined by two independent real parameters: four from the two complex numbers, minus one normalization condition, minus one overall phase. This is exactly the number needed to specify an ordinary spatial direction.

\subsection{Question \& Answer}

If measurements along \(z\) give only discrete values, what becomes of directions between up and down?

They reappear as superpositions. A state need not be an eigenstate of \(L_z\) in order to represent a definite spin direction in space. The simplest examples are
\begin{equation}
\lvert x,+\rangle=\frac{\lvert+\rangle+\lvert-\rangle}{\sqrt2},
\qquad
\lvert y,+\rangle=\frac{\lvert+\rangle+i\lvert-\rangle}{\sqrt2}.
\end{equation}
These are normalized versions of the informal examples given in the lecture.

Consider \(\lvert x,+\rangle\). If we measure \(L_z\), we obtain \(+\hbar/2\) and \(-\hbar/2\) with equal probability \(1/2\). Therefore the average \(z\)-component is
\begin{equation}
\langle L_z\rangle
=
\frac12\left(\frac{\hbar}{2}\right)
+
\frac12\left(-\frac{\hbar}{2}\right)
=0.
\end{equation}
So the actual outcomes remain discrete, but the average behaves as though the spin had been rotated into the \(x\)-direction. This is the lecture's point: the averages behave classically, while individual measurements remain quantized.

In this sense the directions in between are not missing at all. They are encoded in the complex coefficients \(\alpha\) and \(\beta\). To say that the spin points along some chosen axis means that if we measure the component along that axis, we obtain a definite answer. What changes from axis to axis is not the discreteness of measurement outcomes, but which linear combination of \(\lvert+\rangle\) and \(\lvert-\rangle\) represents the aligned state.

The lecture briefly points forward here to helicity, the projection of spin along the direction of motion, but only as a forward marker. The present discussion concerns spin in the rest frame, where the essential issue is already visible in the simpler language of superposition.

\section{Spin Statistics, Composites, And Pauli}

The closing move of the lecture is broader. Spin is not only a label for multiplets; it is correlated with the basic distinction between fermions and bosons. The rule is stated as a theorem of relativistic quantum mechanics:
\begin{theorem}
Particles of integer spin are bosons, and particles of half-integer spin are fermions.
\end{theorem}
Thus spin \(0\), spin \(1\), and spin \(2\) are bosonic, while spin \(\tfrac12\), spin \(\tfrac32\), and so forth are fermionic.

Once this is stated, the lecture immediately turns to composites. Adding two half-integer spins produces an integer. So any object made from an even number of fermionic constituents is a boson, while any object made from an odd number of fermionic constituents is a fermion. The examples are intentionally concrete:
\begin{itemize}
\item a proton and an electron make ordinary hydrogen, which is bosonic as a whole;
\item a proton and a neutron make a deuteron, which is bosonic;
\item a deuteron together with an electron makes deuterium, which is fermionic;
\item an electron and a positron make positronium, which is bosonic.
\end{itemize}
The lecture also remarks that particles and antiparticles have the same spin.

The discussion then returns to the Pauli exclusion principle. The point is subtle but essential: Pauli did not say that two electrons cannot occupy the same orbital motion. He said that no two electrons can occupy the same full quantum state. If two electrons are in the same orbital state, then they must differ in spin. In the familiar atomic setting that means one may be spin up and the other spin down relative to the chosen axis. The exclusion principle acts on the whole state, not on orbital motion alone.

This late atomic aside is still governed by the same angular-momentum counting. In spectroscopic notation, \(S\), \(P\), \(D\), \(F\) correspond to orbital angular momenta \(l=0,1,2,3\). The same rule that gave \(2s+1\) spin states gives \(2l+1\) orbital states: one for \(l=0\), three for \(l=1\), five for \(l=2\), and so on. This is why the angular-momentum formalism developed for spin also reappears in atomic structure.

The lecture closes by deliberately leaving one paradox unresolved. If two fermions can combine into a boson, how can two such composite bosons occupy the same state when their constituents, taken individually, cannot? The lecture explicitly postpones that question. It should remain open here as well; it is a forward pointer, not a solved subsection.

\section{Summary}

The lecture begins with ordinary rotation and ends with particle statistics, but the spine is remarkably tight. Spin is what remains of angular momentum in the center-of-mass rest frame. The angular-momentum operators satisfy a closed algebra, and from that algebra alone one obtains ladder operators, finite multiplets, integer and half-integer families, and the invariant eigenvalue \(L^2=s(s+1)\hbar^2\). For spin \(\tfrac12\), superposition restores the full geometry of orientation even though any chosen component has only discrete outcomes. From there the phenomenological payoff is immediate: integer-spin particles are bosons, half-integer-spin particles are fermions, and Pauli exclusion is a statement about the full quantum state.